% Slide 2 | M2 (Nitin Sureshkumar) | 1:00 | the question and the data
\begin{frame}{The question and the data}
\framesubtitle{Slide 2 \textbullet\ M2 \textbullet\ 1:00 \textbullet\ do two related stocks drift apart for good, or does their gap close?}
\vspace{-2pt}
\begin{columns}[T,onlytextwidth]
\column{0.52\textwidth}
  \footnotesize\raggedright
  \textbf{\color{deckAccent}The idea}
  \begin{itemize}\setlength{\itemsep}{1pt}
    \item two related stocks, such as KO and PEP, tend to move together
    \item when the gap opens, buy the cheap one and sell the expensive one
    \item it pays only if the gap \textbf{really closes}
  \end{itemize}
  \vspace{3pt}
  \textbf{\color{deckAccent}The spread}\\[2pt]
  \centering $X_t=\ln P^A_t-\beta\ln P^B_t-\mu$\\[2pt]
  \raggedright\scriptsize $\beta,\mu$ from a regression on the formation year only
\column{0.44\textwidth}
  \scriptsize\setlength{\tabcolsep}{4pt}
  \begin{tabular}{@{}lll@{}}
    \toprule
    & US & India (NSE) \\
    \midrule
    Stocks         & 56 large caps  & 26 \\
    Sectors        & 7              & 5 \\
    Period         & 2010--2026     & 2012--2026 \\
    Candidate pairs per window & 196 & 55 \\
    Cost per trade, per stock  & 10 bp & 15 bp \\
    \bottomrule
  \end{tabular}\\[4pt]
  \footnotesize
  Daily adjusted closes from Yahoo Finance; pairs only within a sector.
\end{columns}
\pairTakeaway{Today's large caps give the sample \textbf{survivorship bias}, so any edge we find would be overstated.}
\end{frame}
